% chapters/02A_server_linux_docker.tex
\chapter{Server Path A (recommended): Linux + Docker}
\label{ch:server_linux_docker}

\begin{center}
  \vspace{0.6cm}
  \includegraphics[width=0.38\textwidth]{serivdor_pasea_pinguino}
  \vspace{0.8cm}
\end{center}


\section{Goal (what you will achieve)}
This chapter sets up a \textbf{real Linux server foundation} for the rest of the manual.

By the end, you will have:
\begin{itemize}[leftmargin=*]
  \item a Linux machine you can administer (locally or via SSH),
  \item a working Docker installation,
  \item Docker Compose (\texttt{docker compose}) available,
  \item a clean and safe data layout under \texttt{/srv},
  \item basic security hygiene (optional but strongly recommended).
\end{itemize}

\section{Who should choose this path?}
Choose \textbf{Path A} if any of the following is true:
\begin{itemize}[leftmargin=*]
  \item you want an always-on machine (a ``real server''),
  \item you want the most reliable and standard setup,
  \item you prefer fewer layers (no Windows + WSL complexity),
  \item you plan to run multiple services long-term.
\end{itemize}

If you are on Windows and want a guided learning approach first, use Path B (Windows + WSL2).

\section{What is a terminal, and why do we use it?}
A \textbf{terminal} is a text-based interface to your system. You type commands; the system executes them.

Many server tasks are easier and more precise in the terminal:
\begin{itemize}[leftmargin=*]
  \item installing software,
  \item configuring services,
  \item checking logs and system status,
  \item repeating steps reliably (copy/paste).
\end{itemize}

\section{What is \texttt{sudo}?}
\texttt{sudo} means ``run as administrator'' on Linux.
\begin{itemize}[leftmargin=*]
  \item Installing software and changing system configuration requires admin rights.
  \item \texttt{sudo} grants those rights for one command at a time.
\end{itemize}

\section{Step 1 --- Install Linux (high-level, beginner-friendly)}
If Linux is already installed and you can open a terminal, skip to Step 2.

\subsection{Recommended distributions}
For beginners, choose one of these:
\begin{itemize}[leftmargin=*]
  \item \textbf{Ubuntu Server LTS} (very common, lots of tutorials),
  \item \textbf{Debian} (very stable, slightly more conservative).
\end{itemize}

This manual uses commands compatible with both Ubuntu and Debian.

\subsection{During installation: minimal checklist}
When the installer asks questions:
\begin{itemize}[leftmargin=*]
  \item Create a user account (this will be your admin user).
  \item Enable \textbf{OpenSSH server} if offered (recommended).
  \item Choose a hostname (e.g., \texttt{sovereignty-server}).
  \item Keep disk layout simple (we will store our data under \texttt{/srv}).
\end{itemize}

\subsection{After installation: how to open the Linux terminal}
You have two common options:

\paragraph{Option A: local login (keyboard + screen).}
You will see a login prompt. After logging in, you are already in a terminal.

\paragraph{Option B: remote login via SSH (recommended).}
From another computer on the same network, open a terminal (or PowerShell on Windows) and type:
\begin{lstlisting}
ssh username@SERVER_IP
\end{lstlisting}
After login, you are inside the server terminal.

\subsection{Quick check: do you have admin privileges?}
Run:
\begin{lstlisting}
sudo whoami
\end{lstlisting}
If it prints \texttt{root}, your user can run admin commands (good).

\section{Step 2 --- Update Linux (important before installing Docker)}

\encourage{If this is your first time using a terminal: perfect. You do not need to ``understand everything'' today. Just follow the steps and you will learn naturally as you go.}

\subsection{2.1 Refresh package lists}
\begin{lstlisting}
sudo apt update
\end{lstlisting}
\textbf{What this does:} downloads the latest list of available packages. \\
\textbf{Why you do it:} you want current versions and security fixes.

\subsection{2.2 Upgrade installed packages}
\begin{lstlisting}
sudo apt upgrade -y
\end{lstlisting}
\textbf{What this does:} upgrades installed software. \\
\textbf{Why you do it:} reduces bugs and security issues. \\
\textbf{About \texttt{-y}:} automatically answers ``yes'' to prompts.

\subsection{2.3 Install basic tools (recommended)}
\begin{lstlisting}
sudo apt install -y curl ca-certificates gnupg lsb-release unzip
\end{lstlisting}

\textbf{What these are:}
\begin{itemize}[leftmargin=*]
  \item \texttt{curl}: download files from the internet (we use it to fetch Docker keys),
  \item \texttt{ca-certificates}: makes HTTPS downloads verifiable and safe,
  \item \texttt{gnupg}: used to verify signed packages (security),
  \item \texttt{lsb-release}: helps detect your Linux version automatically,
  \item \texttt{unzip}: useful for downloaded archives.
\end{itemize}

\section{Step 3 --- Create a clean data layout}
A common beginner mistake is to store important data \emph{inside} containers.
Do \textbf{not} do that.

\subsection*{The idea (in one sentence)}
Containers are disposable; your data should live in \emph{normal folders on disk}.

\subsection*{Choose a single root folder for persistent data}
Two common options are:
\begin{itemize}[leftmargin=*]
  \item \texttt{/srv} (very common on servers),
  \item \texttt{/opt} (also used for application data).
\end{itemize}

In this manual we use \texttt{/srv}.

\subsection*{Create the folders}
\begin{lstlisting}
sudo mkdir -p /srv/{docker,photos,music,drive,ai,backups}
sudo mkdir -p /srv/docker/compose
\end{lstlisting}

\textbf{What this does:}
\begin{itemize}[leftmargin=*]
  \item \texttt{mkdir} creates directories,
  \item \texttt{-p} means ``create parents if needed'' and ``do not error if it already exists'',
  \item \texttt{\{...\}} creates multiple folders quickly.
\end{itemize}

\textbf{Why you do it:} a consistent structure makes backups, upgrades, and troubleshooting much easier.

\subsection*{Give your user ownership}
\begin{lstlisting}
sudo chown -R $USER:$USER /srv
\end{lstlisting}

\textbf{What this does:}
\begin{itemize}[leftmargin=*]
  \item changes ownership of \texttt{/srv} to your user,
  \item \texttt{-R} applies to everything inside.
\end{itemize}

\textbf{Why you do it:} you can manage configs and volumes without constantly using \texttt{sudo}.

\subsection{Why this matters}
When you upgrade or recreate containers, your data remains safe in \texttt{/srv}.
Containers can be disposable; your files should not be.

\section{Step 4 --- Install Docker Engine (from the official repository)}
The safest approach is to install Docker from Docker's official repository.
This gives you up-to-date packages and security fixes.

\subsection{4.1 Add Docker's official GPG key}
\begin{lstlisting}
sudo install -m 0755 -d /etc/apt/keyrings
curl -fsSL https://download.docker.com/linux/ubuntu/gpg | sudo gpg --dearmor -o /etc/apt/keyrings/docker.gpg
sudo chmod a+r /etc/apt/keyrings/docker.gpg
\end{lstlisting}

\textbf{What this does:}
\begin{itemize}[leftmargin=*]
  \item creates a folder for trusted repository keys,
  \item downloads Docker's signing key and stores it for APT,
  \item makes it readable by the package manager.
\end{itemize}

\textbf{Why you do it:} APT uses the key to verify that packages are authentic (not modified).


\encourage{This step looks scary because it mentions keys and repositories. You are not ``hacking'' anything: you are simply telling Linux to trust the official Docker source. Copy/paste exactly as shown.}

\subsection{4.2 Add the Docker repository (choose Ubuntu or Debian)}
\textbf{Important: pick the correct block for your system.}

\subsubsection*{Ubuntu}
\begin{lstlisting}
echo \
  "deb [arch=$(dpkg --print-architecture) signed-by=/etc/apt/keyrings/docker.gpg] https://download.docker.com/linux/ubuntu \
  $(. /etc/os-release && echo "$VERSION_CODENAME") stable" | \
  sudo tee /etc/apt/sources.list.d/docker.list > /dev/null
\end{lstlisting}

\subsubsection*{Debian}
\begin{lstlisting}
echo \
  "deb [arch=$(dpkg --print-architecture) signed-by=/etc/apt/keyrings/docker.gpg] https://download.docker.com/linux/debian \
  $(. /etc/os-release && echo "$VERSION_CODENAME") stable" | \
  sudo tee /etc/apt/sources.list.d/docker.list > /dev/null
\end{lstlisting}

\textbf{What this does:}
\begin{itemize}[leftmargin=*]
  \item tells APT where to download Docker packages from,
  \item automatically detects your CPU architecture (e.g., amd64),
  \item automatically detects your distro codename (e.g., \texttt{jammy}, \texttt{bookworm}),
  \item restricts trust to packages signed by Docker's key.
\end{itemize}

\textbf{Why you do it:} so \texttt{apt} can install Docker from the official source.

\subsection{4.3 Install Docker Engine and plugins}
\begin{lstlisting}
sudo apt update
sudo apt install -y docker-ce docker-ce-cli containerd.io docker-buildx-plugin docker-compose-plugin
\end{lstlisting}

\textbf{What this installs:}
\begin{itemize}[leftmargin=*]
  \item \texttt{docker-ce}: Docker Engine (Community Edition),
  \item \texttt{docker-ce-cli}: Docker command-line tool,
  \item \texttt{containerd.io}: the container runtime used by Docker,
  \item \texttt{docker-buildx-plugin}: modern image-building features,
  \item \texttt{docker-compose-plugin}: Compose integration (\texttt{docker compose}).
\end{itemize}

\textbf{Why you do it:} you get Docker + Compose in a modern, supported configuration.

\subsection{4.4 Enable and start Docker}
\begin{lstlisting}
sudo systemctl enable docker
sudo systemctl start docker
\end{lstlisting}

\textbf{What this does:}
\begin{itemize}[leftmargin=*]
  \item starts Docker now,
  \item ensures Docker starts automatically after a reboot.
\end{itemize}

\textbf{Why you do it:} your services should come back after restarts.

\section{Step 5 --- Allow running Docker without \texttt{sudo} (recommended)}
By default, Docker commands require \texttt{sudo}. For everyday use, add your user to the \texttt{docker} group:

\begin{lstlisting}
sudo usermod -aG docker $USER
\end{lstlisting}

\textbf{What this does:} allows your user to communicate with the Docker daemon without admin privileges. \\
\textbf{Why you do it:} it makes daily work smoother.

Then log out and log back in (or reboot) so the group membership applies.

\section{Step 6 --- Verify the installation}
\subsection{Docker version}
\begin{lstlisting}
docker --version
\end{lstlisting}
\textbf{Why this matters:} confirms Docker is installed and accessible.

\subsection{Compose version}
\begin{lstlisting}
docker compose version
\end{lstlisting}
\textbf{Why this matters:} confirms Compose is available for multi-service setups.

\subsection{Hello-world test}
\begin{lstlisting}
docker run --rm hello-world
\end{lstlisting}

\textbf{What this does:}
\begin{itemize}[leftmargin=*]
  \item downloads a small test image (if needed),
  \item runs it in a container,
  \item deletes the container after exit (\texttt{--rm}).
\end{itemize}

\textbf{Why this matters:} proves Docker can pull images and run containers correctly.

\encourage{If you saw the hello-world success message: congratulations. You have Docker working. From here on, it is mostly repeating the same pattern.}


\section{Step 7 --- (Optional) Basic firewall and SSH hygiene}
If you only use services inside your home network or behind a VPN, you can keep this simple.
If the server is exposed to the internet, do \textbf{not} skip this.

\subsection{Enable UFW (simple firewall)}
\begin{lstlisting}
sudo apt install -y ufw
sudo ufw default deny incoming
sudo ufw default allow outgoing
sudo ufw allow OpenSSH
sudo ufw enable
sudo ufw status
\end{lstlisting}

\textbf{What this does:}
\begin{itemize}[leftmargin=*]
  \item blocks incoming connections by default,
  \item allows outgoing traffic,
  \item allows SSH so you do not lock yourself out,
  \item enables the firewall and shows status.
\end{itemize}

\textbf{Why this matters:} it reduces your attack surface.

\subsection{SSH: disable password login (recommended)}
Once you can log in via SSH keys, disable password login:

\begin{lstlisting}
sudo nano /etc/ssh/sshd_config
\end{lstlisting}

\textbf{What this does:} opens the SSH server configuration file. \\
\textbf{Why this matters:} password logins can be brute-forced; keys are much safer.

Find the line (or add it) and set:

\begin{lstlisting}
PasswordAuthentication no
\end{lstlisting}

Then restart SSH:

\begin{lstlisting}
sudo systemctl restart ssh
\end{lstlisting}

\textbf{What this does:} applies the new SSH configuration. \\
\textbf{Why this matters:} changes require a restart/reload to take effect.

\section{Verification checklist}
Before moving on, confirm:
\begin{itemize}[leftmargin=*]
  \item \texttt{docker run --rm hello-world} works,
  \item \texttt{docker compose version} prints a version,
  \item \texttt{/srv} exists and contains your data folders,
  \item you understand the core principle: \textbf{containers are disposable; data is not}.
\end{itemize}

\section{Troubleshooting (common issues)}
\subsection{``Permission denied'' when running Docker}
If Docker commands fail without \texttt{sudo}, your user may not yet be in the \texttt{docker} group,
or the group membership has not been applied.

Fix:
\begin{lstlisting}
sudo usermod -aG docker $USER
\end{lstlisting}
Then log out and log back in (or reboot).

\subsection{Docker service not running}
Check status:
\begin{lstlisting}
sudo systemctl status docker
\end{lstlisting}
Start it:
\begin{lstlisting}
sudo systemctl start docker
\end{lstlisting}

\subsection{Compose not found}
Install the Compose plugin:
\begin{lstlisting}
sudo apt install -y docker-compose-plugin
\end{lstlisting}

\section{Next step}
Next, we will deploy the first real service (PhotoPrism) using Docker Compose.
We will reuse the same principles:
\textbf{clear volumes}, \textbf{documented configuration}, and \textbf{easy upgrades}. 